\documentclass[12pt]{article}
\usepackage[T1]{fontenc}
\usepackage[utf8]{inputenc}
\usepackage{lmodern}
\usepackage{textcomp}
\usepackage{enumitem}
\usepackage{geometry}
\geometry{margin=1in}
\begin{document}
\begin{center}
Answer Key - Health Sciences - Graduate Level Assessment 13 \
\small (Answers are ordered exactly as in the question paper.)
\end{center}

\section*{Multiple Choice}
\begin{enumerate}[label=\arabic*.]
\item \textbf{Answer:} C
\\ \textit{Options:} A) Low motor unit strength; B) Low glycogen storage capacity; C) High fatigue resistance; D) Low oxidative capacity; E) Low glycolytic capacity
\\ \textit{Note:} Type IIb muscle fibers are characterized by their high power output and low fatigue resistance, making "high fatigue resistance" an inaccurate descriptor for this fiber type.
\item \textbf{Answer:} D
\\ \textit{Options:} A) Increases in life expectancy at a specified age; B) The increase in the rate of technological advancements; C) The decrease in childhood mortality rates; D) The steadily increasing birth rate; E) The Baby Boom generation
\\ \textit{Note:} The steadily increasing birth rate is not a major reason for the increase in numbers of older adults; rather, the aging population is primarily due to declining mortality rates and the aging of the baby boomer generation.
\item \textbf{Answer:} D
\\ \textit{Options:} A) lipogenesis.; B) lipolysis.; C) lipo-genesis.; D) gluconeogenesis.; E) proteolysis.
\\ \textit{Note:} Gluconeogenesis is the metabolic process through which glucose is synthesized from non-carbohydrate precursors such as lactate, glycerol, or amino acids, making it essential for maintaining blood glucose levels during fasting or intense exercise.
\item \textbf{Answer:} C
\\ \textit{Options:} A) 10-15\%; B) \textgreater{}20\%; C) \textless{}10\%; D) \textgreater{}15\%; E) \textless{}5\%
\\ \textit{Note:} The World Health Organization recommends that free sugars should comprise less than 10\% of total energy intake to reduce the risk of obesity, dental caries, and noncommunicable diseases.
\item \textbf{Answer:} A
\\ \textit{Options:} A) lncRNA; B) piRNA; C) rRNA; D) snRNA; E) tRNA
\\ \textit{Note:} lncRNA, or long non-coding RNA, is indeed a type of non-coding RNA, making it incorrect to categorize it as "not a non-coding RNA."
\end{enumerate}

\section*{True / False}
\begin{enumerate}[label=\arabic*.]
\setcounter{enumi}{5}
\item \textbf{Answer:} False
\\ \textit{Options:} A) True; B) False
\\ \textit{Note:} Proprioception pathways primarily decussate at the level of the spinal cord or brainstem nuclei, but they do not decussate within the corticospinal tract itself, which primarily carries motor signals.
\item \textbf{Answer:} False
\\ \textit{Options:} A) True; B) False
\\ \textit{Note:} The karyotype 47,XX,+13 indicates the presence of an extra chromosome 13, which is associated with Patau syndrome, not Down syndrome, which is characterized by an extra chromosome 21 (47,XX,+21).
\item \textbf{Answer:} False
\\ \textit{Options:} A) True; B) False
\\ \textit{Note:} Medially directed strabismus, or esotropia, is primarily caused by an imbalance in the muscles controlling eye movement, often due to issues with the oculomotor nerve rather than damage to the abducens nerve, which primarily affects lateral eye movement.
\item \textbf{Answer:} True
\\ \textit{Options:} A) True; B) False
\\ \textit{Note:} Humans are known to have a significantly longer lifespan compared to most other animals, with many individuals living beyond 80 years due to advancements in healthcare, nutrition, and living conditions.
\item \textbf{Answer:} False
\\ \textit{Options:} A) True; B) False
\\ \textit{Note:} Senescence is primarily characterized by the cessation of cellular division and growth, often leading to cellular aging and dysfunction rather than continued proliferation.
\end{enumerate}

\section*{Long Form Response}
\begin{enumerate}[label=\arabic*.]
\setcounter{enumi}{10}
\item \textbf{Answer:} Vitamin A plays a critical role in calcium absorption within the small intestine by facilitating the synthesis of calcium-binding proteins, particularly calbindin. Retinoic acid, the active metabolite of vitamin A, enhances the expression of genes responsible for these proteins, thereby increasing the efficiency of calcium uptake from dietary sources. This mechanism underscores the interdependence of vitamins and minerals in nutritional biochemistry, suggesting that adequate vitamin A levels are essential for optimal calcium metabolism. Furthermore, deficiencies in vitamin A can lead to impaired calcium absorption, potentially contributing to bone health issues such as osteoporosis. Thus, ensuring sufficient intake of vitamin A is vital for maintaining overall health and proper skeletal function.
\\ \textit{Note:} correct because it identifies the role of retinoic acid, the active form of vitamin A, in promoting the expression of calcium-binding proteins such as calbindin, which is essential for enhancing calcium absorption in the small intestine and illustrates the interrelationship between vitamins and minerals in nutrition.
\item \textbf{Answer:} The hyoid bone develops from the third and fourth pharyngeal arches during embryogenesis, specifically arising from the mesenchyme of these arches. The third arch contributes to the body of the hyoid, while the fourth arch is responsible for the lesser horns of the bone. This developmental pathway is essential as it influences the hyoid's relationship with adjacent anatomical structures, such as the muscles of the neck and the tongue, which are crucial for functions like swallowing and speech. Furthermore, the hyoid bone serves as an important attachment point for various muscles and ligaments, highlighting the significance of its embryological origins in the broader context of craniofacial development and functionality. Understanding these relationships aids in comprehending various clinical conditions that may arise from anomalies in pharyngeal arch development.
\\ \textit{Note:} correct because it accurately describes the hyoid bone's embryological development from the third and fourth pharyngeal arches and highlights the significance of this development in relation to its anatomical connections and functional roles in the neck, swallowing, and speech.
\end{enumerate}

\vfill
\noindent\textit{End of Answer Key}
\end{document}